 \documentclass[12pt]{article}
\pagestyle{myheadings}

%Enter your name, the portfolio problem number, and the draft number.
\title{Portfolio Problem 1 -- Draft 1}
\author{Ted Sundstrom}

%Enter your name, the portfolio problem number, and the draft number.  This will be a heading on pages after the first page.
\markright{Ted Sundstrom Problem 1 -- Draft 1}


\usepackage{amsmath,amssymb,amsthm,amsfonts,graphics}

%The following commands allow us to typeset theorems, propositions, definitions, etc.
\theoremstyle{plain}
\newtheorem{theorem}{Theorem}
\newtheorem{lemma}[theorem]{Lemma}
\newtheorem{corollary}[theorem]{Corollary}
\newtheorem{proposition}[theorem]{Proposition}
\newtheorem*{definition}{Definition}

\renewcommand{\qedsymbol}{\ensuremath{\blacksquare}}


\begin{document}
\maketitle

%Enter your email address.
\begin{center}
\textbf{email address: sundstrt@gvsu.edu}
\end{center}

\begin{proposition}
If $a$ and $b$ are type 2 integers, then $a \cdot b$ is a type 1 integer.
\end{proposition}

\begin{proof}
We assume that $a$ and $b$ are type 2 integers and will prove that $a \cdot b$ is a type 1 integer.  Since $a$ and $b$ are type 2 integers, there exist integers $m$ and $n$ such that
\[
a = 3m + 2 \text{ ~~~~~~and~~~~~ } b = 3n + 2.
\]
We can now  use substitution and algebra .........
\begin{align*}
ab &= (3m + 2)(3n + 2) \\
    &= 9mn + 6m + 6n + 4 \\
    &= 9mn + 6m + 6n + 3 + 1
\end{align*}
\end{proof}

\end{document}


